\section{Rangkuman}
\label{sec:rangkuman-bab16}

Pada bab ini telah dibahas:

\begin{itemize}
    \item \eng{Disassembly} adalah pekerjaan mengubah kode mesin
          kembali menjadi \eng{mnemonic}. Pekerjaan ini searah dan
          tidak dapat membalikkan \eng{assembly} sepenuhnya, karena
          nama simbolik, komentar, dan batas data hilang saat
          perakitan.

    \item Perintah \code{U} pada DEBUG menampilkan alamat, byte
          mesin, dan \eng{mnemonic} sekaligus. Perkakas lain yang
          sejenis adalah \eng{disassembler} Emu8086, GUI Turbo
          Assembler, \code{ndisasm}, dan \code{objdump}.

    \item Keterbatasan utama \eng{disassembly} adalah sulitnya
          membedakan kode dan data, munculnya alamat relatif dan
          \eng{displacement} menggantikan nama variabel, serta
          hilangnya semua penanda yang berguna bagi manusia.

    \item Pola hasil \eng{disassembly} dapat dikenali: pembandingan
          dan percabangan, perulangan, \eng{prologue} dan
          \eng{epilogue} prosedur, inisialisasi segmen, serta
          pemanggilan interupsi.

    \item Empat metrik optimasi adalah jumlah instruksi yang
          dieksekusi, perkiraan jumlah siklus, ukuran kode dalam
          byte, dan jumlah akses memori. Tanpa pengukuran, optimasi
          hanyalah dugaan yang belum terbukti.

    \item Teknik optimasi level instruksi meliputi: memilih register
          daripada memori, mengganti \asm{MUL}/\asm{DIV} dengan
          pergeseran, memakai \asm{LEA} untuk aritmatika yang tidak
          boleh mengubah \eng{flag}, mengurangi lompatan,
          \eng{loop unrolling}, dan \lstasm|XOR reg, reg| sebagai
          pengganti \lstasm|MOV reg, 0|.

    \item \eng{Strength reduction} menggantikan operasi mahal dengan
          operasi murah, tetapi hanya sahih bila syaratnya dipenuhi,
          misalnya pengali harus pangkat dua dan nilai harus tak
          bertanda.

    \item Optimasi ukuran dan optimasi kecepatan sering bertentangan.
          Pilihan ditentukan oleh batasan proyek, bukan oleh selera.

    \item Pengukuran yang adil menuntut masukan sama, jumlah iterasi
          sama, batas pengukuran yang sama, dan pengulangan beberapa
          kali.

    \item Batas etika optimasi adalah kebenaran lebih dahulu,
          larangan optimasi dini, serta kewajiban menjaga
          keterbacaan dan dokumentasi.
\end{itemize}

\begin{checklistbox}
Tandai setiap butir yang sudah Anda kuasai.

\begin{itemize}
    \item[\checkbox] Saya dapat menjelaskan perbedaan
          \eng{assembly} dan \eng{disassembly} beserta keterbatasan
          \eng{disassembly}.
    \item[\checkbox] Saya dapat membongkar program dengan perintah
          \code{U} dan membaca keluaran alamat, byte, serta
          \eng{mnemonic}.
    \item[\checkbox] Saya dapat mengenali pola perulangan,
          percabangan, \eng{prologue} prosedur, dan pemanggilan
          interupsi pada hasil \eng{disassembly}.
    \item[\checkbox] Saya dapat menghitung jumlah instruksi yang
          dieksekusi, ukuran byte, perkiraan siklus, dan jumlah akses
          memori.
    \item[\checkbox] Saya dapat menerapkan minimal empat teknik
          optimasi level instruksi beserta syarat berlakunya.
    \item[\checkbox] Saya dapat menjelaskan trade-off antara optimasi
          ukuran dan optimasi kecepatan.
    \item[\checkbox] Saya dapat merancang pengukuran yang adil dan
          mencatat waktu dengan \code{INT 1Ah}.
    \item[\checkbox] Saya dapat menjelaskan alasan mengapa optimasi
          dini harus dihindari.
\end{itemize}

\vspace{0.4em}
Bila masih ada butir yang belum bertanda, ulangi bagian yang
bersesuaian sebelum mengerjakan evaluasi.
\end{checklistbox}

\begin{infobox}
\textbf{Rujukan bab.} Materi pada bab ini diperkuat oleh
\cite{wikibooksdisasm}, \cite{bartlett2023}, dan \cite{nasm2024}.
Untuk pembahasan yang lebih mendalam mengenai keterbatasan
\eng{disassembly} dan cara menafsirkan keluaran perkakas pembongkar,
lihat \cite{hyde2010}.
\end{infobox}
